\chapter{Observability and Tracing: OpenTelemetry, Latency, and LLMOps}
\label{chap:observability}

\begin{chaptergoalsbox}
\textbf{What You Will Learn:}
\begin{itemize}[leftmargin=15pt]
  \item Why traditional HTTP 200 monitoring breaks on probabilistic, streaming AI APIs.
  \item OpenTelemetry from first principles: Traces, Spans, \texttt{TraceID}/\texttt{SpanID}, W3C \texttt{traceparent} headers, and \texttt{asyncio} \texttt{contextvars}.
  \item OpenTelemetry Generative AI Semantic Conventions (\texttt{gen\_ai.usage.input\_tokens}, \texttt{gen\_ai.usage.output\_tokens}) and FastAPI instrumentation.
  \item Streaming latency profiling: GPU prefill vs decode, KV-cache prefix reuse, Time to First Token (TTFT), and Inter-Token Latency (ITL).
  \item Token economics, prompt caching pricing ledgers, PII redaction, and tail-based sampling.
\end{itemize}

\textbf{What You Should Already Understand:}
\begin{itemize}[leftmargin=15pt]
  \item Basic Python functions, classes, and HTTP client/server requests.
  \item Chapters~3--6 (RAG pipelines, tool-calling agents, and evaluation metrics).
\end{itemize}

\textbf{Why This Topic Exists:}
In conventional web APIs, an HTTP \texttt{200 OK} response in under a second means the system is healthy. In AI engineering, an HTTP \texttt{200 OK} can mask a 40-second tail latency spike, a \$15 runaway agent loop, a silent vector-index fallback, or a leaked credit card number inside a prompt. Distributed tracing gives engineers X-ray vision across every retrieval call, tool invocation, and token stream.

\textbf{Motivating Practical Problem:}
A legal research assistant's $p99$ response latency (the threshold covering the fastest 99\% of requests while the slowest 1\% exceed it) jumps from 3.2 seconds to 28.5 seconds on a Tuesday afternoon. Without a per-request span waterfall separating vector search, LLM prefill, LLM decode, and SQL tool execution, engineers cannot tell which dependency stalled.
\end{chaptergoalsbox}

\section{Why Black-Box APM Fails on AI Systems}

Traditional Application Performance Monitoring (APM) treats an LLM invocation as one opaque outbound call: \texttt{POST https://api.openai.com/v1/chat/completions -> 200 OK}. That hides three realities:
\begin{itemize}[leftmargin=18pt]
  \item \textbf{HTTP 200 Can Still Be Broken}: A 200 response can contain a safety refusal, invalid JSON, or an empty completion.
  \item \textbf{Wildly Asymmetric Costs}: Two 1.5-second requests can differ in cost by $500\times$ if one hits a 200-token cached prompt (\$0.0003) while the other sends 64,000 uncached tokens (\$0.16).
  \item \textbf{Streaming Perception}: Most AI UIs stream words via \textbf{Server-Sent Events (SSE)}—a standard HTTP response kept open so the server pushes \texttt{data:} chunks as each token is produced. Total request duration misses how long the user stared at a blank screen waiting for the very first word.
\end{itemize}

\section{OpenTelemetry From First Principles}

Before recording AI attributes, understand the four building blocks of \textbf{OpenTelemetry (OTel)}, the industry-standard open telemetry specification:

\begin{enumerate}[leftmargin=18pt]
  \item \textbf{Trace and \texttt{TraceID}}: A \textbf{Trace} represents the complete end-to-end journey of one user request across services, identified globally by a 128-bit hexadecimal \texttt{TraceID} (e.g., \texttt{4bf92f3577b34da6a3ce929d0e0e4736}).
  \item \textbf{Span and \texttt{SpanID}}: A \textbf{Span} represents one named, timed operation inside that trace (such as \texttt{VectorSearch}, \texttt{chat.completion}, or \texttt{ToolExecution}). Each span carries its own 64-bit \texttt{SpanID}, its \texttt{ParentSpanID}, start/end timestamps, structured key-value \textbf{attributes}, and timestamped \textbf{events} (like \texttt{"first\_token"}). Plotting parent and child spans by start time produces a \textbf{waterfall diagram} (Figure~\ref{fig:v2-ch07-spans}).
  \item \textbf{W3C \texttt{traceparent} Header}: When Service A calls Service B over HTTP, how does Service B know which trace and parent span it belongs to? Service A injects the W3C TraceContext HTTP header:
  \begin{center}\small\ttfamily
  traceparent: 00-4bf92f3577b34da6a3ce929d0e0e4736-00f067aa0ba902b7-01
  \end{center}
  Its four hyphenated fields are: \texttt{version} (\texttt{00}), \texttt{trace-id} (32 hex chars), \texttt{parent-span-id} (16 hex chars), and \texttt{trace-flags} (\texttt{01} = sampled).
  \item \textbf{Async Context Propagation (\texttt{contextvars})}: Inside a Python \texttt{asyncio} server, hundreds of user requests interleave concurrently on a single thread. A global variable \texttt{CURRENT\_SPAN} would scramble traces across users! Instead, Python provides \texttt{contextvars.ContextVar}—a task-local slot that Python's event loop automatically swaps whenever it switches between coroutines. OpenTelemetry stores the active span inside a \texttt{ContextVar}. \textbf{Watch out}: if you offload work to a separate OS thread (\texttt{loop.run\_in\_executor}) or a background worker queue, a new thread starts with an \emph{empty} \texttt{ContextVar}, breaking the parent-child link unless you explicitly capture \texttt{ctx = otel\_context.get\_current()} and re-attach it inside the worker.
\end{enumerate}

\begin{figure}[htbp]
\centering
\begin{tikzpicture}[
  root/.style={draw=headingblue, fill=codebg, rounded corners=2pt, font=\sffamily\scriptsize\bfseries, minimum width=9.6cm, minimum height=0.8cm},
  agent/.style={draw=exampleBlue, fill=exampleBg, rounded corners=2pt, font=\sffamily\scriptsize, minimum width=8.8cm, minimum height=0.8cm},
  child/.style={draw=intuitionTeal, fill=intuitionBg, rounded corners=2pt, font=\sffamily\tiny\bfseries, minimum width=2.8cm, minimum height=0.75cm, align=center},
  arrow/.style={-{Stealth[scale=0.8]}, thick, color=headingblue}
]

\node[root] (session) at (0, 3.2) {Trace \texttt{4bf92f35...}: \texttt{POST /chat} (Root Span)};
\node[agent] (chain) at (0, 2.0) {Child Span: \texttt{AgentWorkflow.execute} (\texttt{gen\_ai.system="langgraph"})};

\node[child] (s1) at (-3.1, 0.6) {Span: \texttt{VectorSearch}\\\texttt{db.system="pgvector"}};
\node[child] (s2) at (0, 0.6) {Span: \texttt{chat.completion}\\\texttt{gen\_ai.request.model="gpt-4o"}};
\node[child] (s3) at (3.1, 0.6) {Span: \texttt{ToolExecution}\\\texttt{tool.name="run\_sql"}};

\draw[arrow] (session.south) -- (chain.north);
\draw[arrow] (chain.south) -- (s1.north);
\draw[arrow] (chain.south) -- (s2.north);
\draw[arrow] (chain.south) -- (s3.north);

\end{tikzpicture}
\caption{OpenTelemetry Span Hierarchy for an AI Agent Request.}
\label{fig:v2-ch07-spans}
\end{figure}

\subsection{OpenTelemetry Semantic Conventions for Generative AI}
Standardizing attribute names via OTel GenAI Semantic Conventions lets you switch backends (Langfuse, Grafana Tempo, Datadog, Honeycomb) without rewriting application instrumentation (Table~\ref{tab:v2-ch07-otel}).

\begin{table}[htbp]
\centering
\small
\begin{tabular}{lll}
\toprule
\textbf{Attribute Name} & \textbf{Type} & \textbf{Description / Example} \\
\midrule
\texttt{gen\_ai.system} & string & Provider identifier (\texttt{openai}, \texttt{anthropic}, \texttt{cohere}) \\
\texttt{gen\_ai.operation.name} & string & Operation (\texttt{chat}, \texttt{embeddings}, \texttt{execute\_tool}) \\
\texttt{gen\_ai.request.model} & string & Model name (\texttt{gpt-4o}, \texttt{claude-3-5-sonnet-20241022}) \\
\texttt{gen\_ai.request.temperature} & float & Sampling temperature ($0.0 \le T \le 2.0$) \\
\texttt{gen\_ai.request.max\_tokens} & int & Maximum output token limit \\
\texttt{gen\_ai.response.id} & string & Provider completion ID (\texttt{chatcmpl-9A8X...}) \\
\texttt{gen\_ai.usage.input\_tokens} & int & Input prompt tokens consumed \\
\texttt{gen\_ai.usage.output\_tokens} & int & Output completion tokens generated \\
\bottomrule
\end{tabular}
\caption{OpenTelemetry Generative AI Semantic Attributes.}
\label{tab:v2-ch07-otel}
\end{table}

Listing~\ref{lst:v2-ch07-fastapi} shows real \texttt{opentelemetry-sdk} instrumentation inside a FastAPI endpoint, including child spans, span events for TTFT, and explicit context propagation across a worker thread.

\begin{lstlisting}[language=Python, label={lst:v2-ch07-fastapi}, caption={FastAPI OpenTelemetry Instrumentation with Context Propagation}]
import asyncio
import time
from fastapi import FastAPI
from opentelemetry import context as otel_context, trace
from opentelemetry.sdk.trace import TracerProvider
from opentelemetry.sdk.trace.export import BatchSpanProcessor, ConsoleSpanExporter

provider = TracerProvider()
provider.add_span_processor(BatchSpanProcessor(ConsoleSpanExporter()))
trace.set_tracer_provider(provider)
tracer = trace.get_tracer("ai.service")

app = FastAPI()

def blocking_vector_lookup(query: str, parent_ctx) -> list[str]:
    # Re-attach asyncio parent context inside worker thread!
    token = otel_context.attach(parent_ctx)
    try:
        with tracer.start_as_current_span("VectorSearch") as span:
            span.set_attribute("db.system", "pgvector")
            time.sleep(0.02)
            return ["doc_42", "doc_99"]
    finally:
        otel_context.detach(token)

@app.post("/chat")
async def chat_endpoint(query: str) -> dict:
    with tracer.start_as_current_span("chat_request") as root:
        ctx = otel_context.get_current()
        docs = await asyncio.to_thread(blocking_vector_lookup, query, ctx)

        with tracer.start_as_current_span("chat.completion") as llm_span:
            llm_span.set_attribute("gen_ai.system", "openai")
            llm_span.set_attribute("gen_ai.request.model", "gpt-4o")
            await asyncio.sleep(0.05)
            llm_span.add_event("first_token", {"ttft_ms": 50.0})
            llm_span.set_attribute("gen_ai.usage.input_tokens", 1240)
            llm_span.set_attribute("gen_ai.usage.output_tokens", 85)
        return {"docs": len(docs), "status": "ok"}
\end{lstlisting}

\textbf{Step-by-Step Walkthrough:}
\begin{enumerate}[leftmargin=18pt]
  \item Lines 8--11 configure \texttt{TracerProvider} with a non-blocking \texttt{BatchSpanProcessor} so exporting telemetry never stalls FastAPI event-loop handlers.
  \item Line 28 opens the parent \texttt{chat\_request} span and captures the active \texttt{ContextVar} via \texttt{otel\_context.get\_current()} on line 29.
  \item Line 30 offloads \texttt{blocking\_vector\_lookup} to a thread pool; lines 17--23 re-attach \texttt{parent\_ctx} so \texttt{VectorSearch} appears as a child of \texttt{chat\_request} rather than splitting into an orphaned trace.
  \item Lines 32--38 record standard \texttt{gen\_ai.*} attributes and emit a \texttt{"first\_token"} event at TTFT.
\end{enumerate}

\section{Streaming Latency: Prefill, Decode, and KV-Cache Reuse}

When an LLM streams tokens over SSE, total duration splits into two physically distinct GPU phases (Figure~\ref{fig:v2-ch07-ttft}):

\begin{quote}
\textbf{GPU Prefill vs Decode \& KV-Cache Prefix Reuse:} During \textbf{Prefill}, the GPU reads all prompt tokens in parallel in one large matrix pass, computing intermediate Key and Value attention states (\textbf{KV cache}) and storing them in high-bandwidth GPU memory (HBM). This step is \emph{compute-bound} and determines \textbf{Time to First Token (TTFT)}. Once the first token appears, the GPU enters \textbf{Decode}: it emits one token at a time, and on every step it must stream the model's weights plus the entire KV cache from HBM into on-chip compute memory (SRAM), making decode \emph{memory-bandwidth bound}. Crucially, if consecutive requests share an identical static prompt prefix (like a 4,000-token system prompt), inference servers can reuse cached KV-cache blocks directly (\textbf{Prompt Caching}), skipping redundant prefill compute and slashing both TTFT and input token fees.
\end{quote}

\begin{figure}[htbp]
\centering
\begin{tikzpicture}[
  bar/.style={draw=headingblue, fill=codebg, rounded corners=2pt, font=\sffamily\scriptsize, align=center, minimum height=0.95cm}
]

\draw[->, thick, color=gray!70] (0,0) -- (12.5,0) node[right, font=\tiny\sffamily] {Time (ms)};
\draw[shift={(0,0)}] (0pt,3pt) -- (0pt,-3pt) node[below, font=\tiny\sffamily] {0 ms (Request)};
\draw[shift={(4.5,0)}] (0pt,3pt) -- (0pt,-3pt) node[below, font=\tiny\sffamily] {450 ms (Token 1)};
\draw[shift={(11.5,0)}] (0pt,3pt) -- (0pt,-3pt) node[below, font=\tiny\sffamily] {1,850 ms (Token 70)};

\node[bar, fill=warningBg, minimum width=4.5cm] at (2.25, 1.25) {
  \textbf{Prefill Phase (TTFT = 450 ms)}\\
  Network RTT + Parallel Prompt Pass
};

\node[bar, fill=intuitionBg, minimum width=7.0cm] at (8.0, 1.25) {
  \textbf{Decode Phase ($1{,}850 - 450 = 1{,}400\text{ ms}$)}\\
  $N = 70$ tokens ($69$ steps) $\implies$ ITL $= \frac{1400}{69} \approx 20.29\text{ ms/tok}$
};

\end{tikzpicture}
\caption{Streaming Latency Decomposition: Token 1 completes at TTFT (450 ms); the remaining $N - 1 = 69$ intervals govern average Inter-Token Latency (ITL).}
\label{fig:v2-ch07-ttft}
\end{figure}

Notice why we divide by $N_{\text{completion}} - 1$ when computing average \textbf{Inter-Token Latency (ITL)}: Token 1 arrives at TTFT, leaving $N - 1$ inter-token intervals across $(t_{\text{end}} - \text{TTFT})$:
\begin{equation}
\text{ITL} = \frac{t_{\text{completion}} - \text{TTFT}}{N_{\text{completion}} - 1}
\end{equation}
Listing~\ref{lst:v2-ch07-stream} tracks TTFT and per-token ITLs on any streaming generator.

\begin{lstlisting}[language=Python, label={lst:v2-ch07-stream}, caption={Streaming TTFT and ITL Latency Profiler}]
import time

class StreamingMetricsCollector:
    def __init__(self):
        self.start_time: float = 0.0
        self.first_token_time: float = 0.0
        self.end_time: float = 0.0
        self.token_timestamps: list[float] = []

    def start(self) -> None:
        self.start_time = time.perf_counter()

    def record_chunk(self) -> None:
        now = time.perf_counter()
        if not self.token_timestamps:
            self.first_token_time = now
        self.token_timestamps.append(now)

    def finalize(self) -> dict[str, float | int]:
        self.end_time = time.perf_counter()
        total_tokens = len(self.token_timestamps)
        ttft = (self.first_token_time - self.start_time) if total_tokens > 0 else 0.0
        total_duration = self.end_time - self.start_time

        itls = [
            self.token_timestamps[i] - self.token_timestamps[i - 1]
            for i in range(1, total_tokens)
        ]
        avg_itl = (sum(itls) / len(itls)) if itls else 0.0
        p95_itl = sorted(itls)[int(len(itls) * 0.95)] if len(itls) >= 20 else avg_itl

        return {
            "total_duration_sec": total_duration,
            "ttft_sec": ttft,
            "total_completion_tokens": total_tokens,
            "avg_itl_ms": avg_itl * 1000.0,
            "p95_itl_ms": p95_itl * 1000.0,
        }
\end{lstlisting}

\textbf{Step-by-Step Trace:} If \texttt{start()} records $t=0.000$ s and 4 tokens arrive at $0.450, 0.470, 0.495, 0.510$ s: \texttt{ttft = 0.450} s, \texttt{itls = [0.020, 0.025, 0.015]} ($N-1 = 3$ intervals), \texttt{avg\_itl\_ms = 20.0} ms, and \texttt{total\_duration\_sec = 0.510} s.

\section{Token Economics and Multi-Tenant Cost Attribution}

Production AI platforms track spend per tenant and feature using three-tier pricing (uncached input, cached prefix input, and output completion tokens):
\begin{equation}
\text{Cost} = \frac{N_{\text{uncached}}}{10^6} P_{\text{in}} + \frac{N_{\text{cached}}}{10^6} P_{\text{cache}} + \frac{N_{\text{out}}}{10^6} P_{\text{out}}
\end{equation}
Table~\ref{tab:model_pricing} lists reference pricing tiers used in Listing~\ref{lst:v2-ch07-cost} and Exercise~\thechapter.2.

\begin{table}[htbp]
\centering
\small
\begin{tabular}{lccc}
\toprule
\textbf{Model Identifier} & \textbf{Input (\$/1M)} & \textbf{Cached Input (\$/1M)} & \textbf{Output (\$/1M)} \\
\midrule
\texttt{gpt-4o} & \$2.50 & \$1.25 & \$10.00 \\
\texttt{claude-3-5-sonnet-20241022} & \$3.00 & \$0.30 & \$15.00 \\
\texttt{gpt-4o-mini} & \$0.15 & \$0.075 & \$0.60 \\
\bottomrule
\end{tabular}
\caption{Illustrative LLM Token Pricing Tiers (USD per 1 Million Tokens).}
\label{tab:model_pricing}
\end{table}

\begin{lstlisting}[language=Python, label={lst:v2-ch07-cost}, caption={Multi-Tenant Token Cost Ledger}]
from dataclasses import dataclass

@dataclass(frozen=True)
class ModelPricing:
    input_per_million: float
    cached_input_per_million: float
    output_per_million: float

PRICING_TABLE = {
    "gpt-4o": ModelPricing(2.50, 1.25, 10.00),
    "claude-3-5-sonnet-20241022": ModelPricing(3.00, 0.30, 15.00),
    "gpt-4o-mini": ModelPricing(0.15, 0.075, 0.60),
}

class CostLedger:
    def __init__(self):
        self.tenant_spend: dict[str, float] = {}

    def record_usage(
        self, tenant_id: str, model: str,
        prompt_tokens: int, cached_tokens: int, completion_tokens: int
    ) -> float:
        pricing = PRICING_TABLE.get(model, PRICING_TABLE["gpt-4o"])
        uncached_prompt = max(0, prompt_tokens - cached_tokens)
        cost = (
            (uncached_prompt / 1_000_000.0) * pricing.input_per_million
            + (cached_tokens / 1_000_000.0) * pricing.cached_input_per_million
            + (completion_tokens / 1_000_000.0) * pricing.output_per_million
        )
        self.tenant_spend[tenant_id] = self.tenant_spend.get(tenant_id, 0.0) + cost
        return cost
\end{lstlisting}

\section{PII Redaction and Tail-Based Sampling}

Exporting raw prompts to third-party trace platforms risks violating GDPR, HIPAA, or PCI-DSS if users paste SSNs, card numbers, or API keys. Scrub attributes in-process before spans leave the service (Figure~\ref{fig:v2-ch07-pipeline} and Listing~\ref{lst:v2-ch07-pii}).

\begin{figure}[htbp]
\centering
\begin{tikzpicture}[
  box/.style={draw=headingblue, fill=codebg, rounded corners=2pt, font=\sffamily\scriptsize, align=center, minimum height=1.0cm},
  arrow/.style={-{Stealth[scale=0.8]}, thick, color=headingblue}
]

\node[box, minimum width=2.7cm] (app) at (0, 0) {AI Service\\(FastAPI / LangGraph)};
\node[box, fill=warningBg, minimum width=3.1cm] (filter) at (3.9, 0) {In-Process PII\\Scrubber};
\node[box, fill=intuitionBg, minimum width=3.2cm] (collector) at (8.0, 0) {OTel Collector\\(Tail Sampling)};
\node[box, fill=exampleBg, minimum width=2.4cm] (langfuse) at (11.9, 0) {Langfuse /\\Grafana};

\draw[arrow] (app) -- (filter) node[midway, above, font=\tiny\sffamily] {Raw Span};
\draw[arrow] (filter) -- (collector) node[midway, above, font=\tiny\sffamily] {Redacted};
\draw[arrow] (collector) -- (langfuse) node[midway, above, font=\tiny\sffamily] {Filtered};

\end{tikzpicture}
\caption{Production Telemetry Pipeline: Scrub PII inside the service process before tail-based sampling at the OTel Collector.}
\label{fig:v2-ch07-pipeline}
\end{figure}

\begin{lstlisting}[language=Python, label={lst:v2-ch07-pii}, caption={In-Process Regex PII Scrubber}]
import re

class PIIScrubber:
    def __init__(self):
        self.patterns = [
            (re.compile(r"\b\d{3}-\d{2}-\d{4}\b"), "[REDACTED_SSN]"),
            (re.compile(r"\b(?:\d{4}[ -]?){3}\d{4}\b"), "[REDACTED_CARD]"),
            (re.compile(r"(sk-[a-zA-Z0-9]{20,}|ghp_[a-zA-Z0-9]{36})"), "[REDACTED_KEY]"),
        ]

    def scrub(self, text: str) -> str:
        for pattern, replacement in self.patterns:
            text = pattern.sub(replacement, text)
        return text
\end{lstlisting}

\begin{underthehoodbox}[Head-Based vs Tail-Based Sampling]
At 10 million daily requests, storing 100\% of verbose prompt spans is prohibitively expensive. \textbf{Head-based sampling} flips a coin at the start of a request (keeping say 1\%)—meaning 99\% of rare errors and slow outliers vanish! \textbf{Tail-based sampling} buffers spans briefly in the OTel Collector until the trace finishes, keeping 100\% of errors and slow traces ($>5$ s) while retaining only 0.5\% of routine fast traces.
\end{underthehoodbox}

\begin{realworldbox}[Isolating a Cascading SQL Tool Bottleneck via Span Waterfall]
On a routine Tuesday afternoon, an enterprise order-assistant's $p99$ latency climbed from 2.1 s to 58 s. Inspecting a single slow trace waterfall in Langfuse pinpointed the culprit immediately:
\begin{lstlisting}[basicstyle=\ttfamily\scriptsize]
TraceID: 91c8e2... (total=58,140 ms)
  |-- VectorSearch (pgvector)          18 ms
  |-- chat.completion (gpt-4o, TTFT)  350 ms
  |-- ToolExecution (InventoryCheck) 56,910 ms  <-- SQL Seq Scan on 40M rows!
  +-- chat.completion (final synth)   862 ms
\end{lstlisting}
The LLM and vector store were healthy; the \texttt{InventoryCheck} tool was running an unindexed \texttt{WHERE sku = ... AND warehouse\_id = ...} query across 40 million rows. Adding a composite B-tree index restored $p99$ latency to 1.8 s within 10 minutes.
\end{realworldbox}

\begin{tryityourselfbox}[Instrument an Async Streaming Route with OTel \& TTFT]
Extend Listing~\ref{lst:v2-ch07-fastapi} so that a simulated streaming helper yields 5 token chunks spaced 30 ms apart after an initial 120 ms prefill delay, records TTFT and average ITL via \texttt{StreamingMetricsCollector}, sets \texttt{gen\_ai.usage.output\_tokens = 5}, and prints the exported span JSON to stdout. (Complete solution in Section~7.8.)
\end{tryityourselfbox}

\section{Chapter Summary \protect\& Key Takeaways}
\begin{itemize}[leftmargin=18pt]
  \item \textbf{OpenTelemetry Core}: Traces link spans across services via 128-bit \texttt{TraceID}s, W3C \texttt{traceparent} headers, and Python \texttt{contextvars} (explicitly attached when crossing threads).
  \item \textbf{GenAI Conventions}: Standardize on \texttt{gen\_ai.usage.input\_tokens} and \texttt{gen\_ai.usage.output\_tokens}.
  \item \textbf{Prefill vs Decode}: TTFT reflects compute-bound prefill (accelerated by KV-cache prefix reuse); ITL reflects memory-bandwidth-bound decode divided across $N - 1$ token intervals.
  \item \textbf{Cost \& Privacy}: Attribute spend per tenant using three-tier pricing, scrub PII before export, and use tail-based sampling so errors are never dropped.
\end{itemize}

\section{Exercises}

\subsection*{Mechanics}
\textbf{Exercise \thechapter.1}: A streaming LLM request takes 2,400 ms total to emit 80 completion tokens, with Token 1 arriving at 400 ms. Compute: (1) TTFT, (2) average ITL in ms/token, and (3) decode throughput in tokens/second.

\textbf{Exercise \thechapter.2}: A service sends 100,000 daily requests to \texttt{claude-3-5-sonnet-20241022}. Each request has a static 4,000-token system prompt, a 500-token dynamic user query, and 300 completion tokens. Compute daily cost without vs with prompt caching using the rates in Table~\ref{tab:model_pricing}.

\subsection*{Debugging}
\textbf{Exercise \thechapter.3}: A team instruments FastAPI with OpenTelemetry, yet \texttt{AgentRun} and \texttt{ToolExecution} show up as completely separate top-level traces with different \texttt{TraceID}s whenever \texttt{ToolExecution} runs in \texttt{loop.run\_in\_executor()}. Diagnose the context failure and provide the fix.

\subsection*{Architecture}
\textbf{Exercise \thechapter.4}: Design a PCI-DSS compliant LLM observability architecture across 50 microservices that scrubs cardholder numbers before network export, tracks per-tenant token spend, retains 100\% of errors, and alerts on daily spend anomalies.

\section{Detailed Solutions}

\subsection*{Solution \thechapter.1}
1. $\text{TTFT} = 400\text{ ms}$. 2. Decode duration $= 2400 - 400 = 2000\text{ ms}$ across $80 - 1 = 79$ intervals, so $\text{ITL} = 2000 / 79 \approx 25.32\text{ ms/token}$. 3. Decode throughput $= 79 / 2.0\text{ s} = 39.5\text{ tokens/sec}$.

\subsection*{Solution \thechapter.2}
From Table~\ref{tab:model_pricing} for \texttt{claude-3-5-sonnet-20241022}: \$3.00/1M input, \$0.30/1M cached input, \$15.00/1M output.
\begin{itemize}
  \item \textbf{Without caching}: Input $= 10^5 \times \frac{4500}{10^6} \times \$3.00 = \$1{,}350.00$; Output $= 10^5 \times \frac{300}{10^6} \times \$15.00 = \$450.00$; Total $= \mathbf{\$1{,}800.00/\text{day}}$.
  \item \textbf{With caching}: Cached prefix $= 10^5 \times \frac{4000}{10^6} \times \$0.30 = \$120.00$; Uncached query $= 10^5 \times \frac{500}{10^6} \times \$3.00 = \$150.00$; Output $= \$450.00$; Total $= \mathbf{\$720.00/\text{day}}$ (60\% savings).
\end{itemize}

\subsection*{Solution \thechapter.3}
OpenTelemetry tracks the active span in an \texttt{asyncio} \texttt{ContextVar}. Worker threads spawned by \texttt{run\_in\_executor()} do not automatically inherit the coroutine's \texttt{ContextVar}, so the child span sees no parent and mints a fresh \texttt{TraceID}. Fix: capture \texttt{ctx = otel\_context.get\_current()} on the event-loop side, pass \texttt{ctx} into the executor function, and wrap the body in \texttt{token = otel\_context.attach(ctx)} / \texttt{otel\_context.detach(token)}.

\subsection*{Solution \thechapter.4}
1. \textbf{In-Process Scrubbing}: Run \texttt{PIIScrubber} inside every microservice before span attributes are serialized. 2. \textbf{Internal Collector Gateway}: Export spans over mTLS gRPC to an internal VPC OTel Collector cluster. 3. \textbf{Tail-Based Sampling}: Retain 100\% of errors and $>5$ s traces, 0.5\% of normal traces, while extracting exact \texttt{gen\_ai.usage.*} metrics into Prometheus/ClickHouse for 100\% accurate tenant billing and anomaly alerts.

\subsection*{Solution (Try It Yourself)}
\begin{lstlisting}[language=Python]
async def stream_with_telemetry():
    collector = StreamingMetricsCollector()
    with tracer.start_as_current_span("chat.completion") as span:
        span.set_attribute("gen_ai.request.model", "gpt-4o")
        collector.start()
        await asyncio.sleep(0.12)  # Simulated prefill
        for idx in range(5):
            collector.record_chunk()
            if idx == 0:
                span.add_event("first_token")
            await asyncio.sleep(0.03)  # Simulated decode step
        metrics = collector.finalize()
        span.set_attribute("gen_ai.usage.output_tokens", metrics["total_completion_tokens"])
        span.set_attribute("llm.ttft_ms", round(metrics["ttft_sec"] * 1000, 2))
        span.set_attribute("llm.avg_itl_ms", round(metrics["avg_itl_ms"], 2))
    provider.force_flush()
\end{lstlisting}

\begin{labbridgebox}[Project Bridge]
Combine evaluation and telemetry by inspecting how \texttt{projects/p07-llm-eval-pipeline/evaluator/harness.py} aggregates per-case latency and scores, and consider where OpenTelemetry spans attach around graph nodes in \texttt{projects/p05-langgraph-agent-state/agent/graph.py}.
\end{labbridgebox}

\section{Further Reading \protect\& Citations}
\begin{itemize}[leftmargin=18pt]
  \item OpenTelemetry. (2025). \emph{Semantic Conventions for Generative AI Systems}. \url{https://opentelemetry.io/docs/specs/semconv/gen-ai/}.
  \item W3C. (2021). \emph{Trace Context Recommendation (\texttt{traceparent})}. \url{https://www.w3.org/TR/trace-context/}.
  \item Langfuse. (2025). \emph{Open Source LLM Engineering Platform Documentation}. \url{https://langfuse.com/docs}.
\end{itemize}
